\scanpage{17}
\handwritten{\mspage{8}}
En caractéristique quelconque, on ne peut pour l'instant que donner des cas assez particuliers ou la conjecture soit vérifiée; par Katz, cela marche en tous cas pour \deleted{les}\inserted{\handwritten{$X$}} intersection\deleted{s} complète\deleted{s}\inserted{\handwritten{et \deleted{\unsure{?}} un pinceau ``en général''}} (ce qui suffit pour faire marcher l'exemple minimal de Griffiths: $X$ quadrique dans $P^5$ ..). \deleted{La méthode} \deleted{p}Pour vérifier que $E^{2n-1}(Y_{\bar s})$ n'est pas de niveau 1, on constate aisément qu'il suffit de vérifier qu'il existe \uline{un} point $t$ du pinceau tel que $E^{2n-1}(Y_{\bar t})$ ne soit pas de niveau 1 (il faut utiliser pour ceci le critère de Néron-Ogg-Chafarévitch de bonne réduction des VA; la triste vérité est que je ne sais pas \deleted{montrer}\inserted{montrer de façon générale}, même en \deleted{\unsure{termes|fermes}} admettant les conjectures standard, que le niveau d'un motif ne peut que diminuer par spécialisation \ldots) Quitte à écrire $k$ comme limite d'anneaux de type fini au sens absolu, la même observation est d'ailleurs valable, en se spécialisant sur un corps fini. Or dans ce cas, le niveau 1 implique\inserted{\handwritten{(et est \unsure{peut-être équivalent à})}} la divisibilité \inserted{\handwritten{par ${\unsure{q|1}}^{n-1}$}} des valeurs propres de frobenius opérant sur $E^{2n-1}(Y_{\bar t})$ \deleted{par $q^{n-1}$} ($q=\mathrm{card}\ k(t)$), et il suffit donc de trouver $t$ tel que l'une des valeurs propres de frobenius soit une unité $p$-adique. C'est ce que Katz vérifie dans divers cas par une utilisation astucieuse de La formule de Lefschetz-Woodshole-Verdier.

7.\quad Le procédé de Katz. A lui de jouer ....
% archive sheet number: 16
